\scanpage{31}
\margin{Car si $g\in\unsure{\mathcal Z|Z|\mathfrak z}(\gamma)$, $g.D=D'$ est conjuguée de $D$ par élément de $T$, donc mod \unsure{ce|cet} dernier $g$ est dans $\mathcal Z(\gamma)\cap T'$}
\diagram{\resizebox{\linewidth}{!}{\begin{tikzcd}[ampersand replacement=\&,row sep=1.4em,column sep=0.8em,row sep=1.9em,column sep=2.3em]
T\arrow[dash,to=1-2,"{\mathcal Z}"]\arrow[dash,to=2-1] \& \mathcal Z(\gamma)\arrow[dash,to=1-3]\arrow[dash,to=2-2] \& N(T)\arrow[dash,to=1-4]\arrow[dash,to=2-3] \& G\arrow[dash,to=2-4] \\
\gamma\in T\cap N(T')\arrow[dash,to=2-2,"{\mathcal Z}"]\arrow[dash,to=3-1,"{\mathbb Z/h\mathbb Z}"] \& \mathcal Z(\gamma)\cap N(T')={N'}^{\gamma}\arrow[dash,to=2-3]\arrow[dash,to=3-2,"{\mathbb Z/h\mathbb Z}"] \& N(T)\cap N(T')\arrow[dash,to=2-4]\arrow[dash,to=3-3] \& N(T')\arrow[dash,to=3-4,"{W'}"] \\
T\cap T'\arrow[dash,to=3-2,"{\mathcal Z}"]\arrow[equals,to=4-1] \& \mathcal Z(\gamma)\cap T'={T'}^\gamma\arrow[equals,to=3-3]\arrow[phantom,to=4-2,"{\rotatebox{90}{$\simeq$}}"] \& N(T)\cap T'\arrow[dash,to=3-4] \& T' \\
0 \& P'/(1-c)P'(1)\arrow[phantom,to=5-2,"{\rotatebox{90}{$\simeq$}}"] \&  \&  \\
 \& \mathcal Z \&  \& 
\end{tikzcd}}}
\margin{Car \unsure{tt|tout} élément de $W$ qui centralise $c$ est une puissance de $c$}
$\uparrow$ Car tout élément de $T$ qui centralise $D$ est \deleted{\unsure{dans}} l'élément unité
\margin{\uline{\unsure{À|Â} prouver}}

Tout élément de $T$ qui normalise $D$ est dans $\unsure{\kappa|k|K}(\mu_h)$, i.e. $h\psi$ (ou, au choix, les $h\alpha_i$) sont \unsure{comb. lin$^{\mathrm{r}}$|comb. lin.|complétées} à coeff. entiers des éléments $\alpha_i+\psi$.

\unsure{Provient du|Revient donc|Prouvons donc} lemme
\[
\det\begin{pmatrix}
n_1+1 & n_2, & n_3 & \cdots & ,\,n_r\\
\unsure{n_2|n_1|n\,x} & n_2+1 & n_3 & \cdots & n_r\\
\cdots & & & &\\
n_r & n_r & & \cdots & n_r+1
\end{pmatrix}=\sum n_i+1
\]
Un élément de $N(T)\cap N(D)$ \deleted{centralise}\inserted{normalise} $D$, et normalise $T$, donc il normalise \deleted{$\unsure{\mu}$} $\mu_h\simeq\mathbb Z/h\mathbb Z\subset T$ qui est formé des éléments de $T$ qui normalisent $D$. Donc
\[
N(T)\cap N(D)=\deleted{T'\cap\mathrm{Norm}(\mu)}\inserted{N(\mu)\cap N(D)}
\quad\bigl(\underset{\substack{\rotatebox{90}{$\simeq$}\\\mathbb Z/h\mathbb Z=\mu_h}}{\mu}\subset T\text{ engendré par }\gamma\bigr)
\]
% archive sheet number: 30
